\section{The AlphaZero algorithm}
\label{sec:alphazero}

\subsection{Three papers, and what each one removed}
\label{sec:three-papers}

The AlphaZero design did not appear at once. It is the residue of three papers,
each of which deleted components that the field believed were necessary, and each
deletion made the system both simpler and stronger. Table~\ref{tab:removals}
records the sequence; the direction of travel is the message.

\begin{table}[t]
\centering
\small
\caption{What each paper removed. Every entry in the first column is a component
that a specialist would have defended before the following paper appeared.}
\label{tab:removals}
\begin{tabularx}{\textwidth}{@{}l l l l@{}}
\toprule
& AlphaGo (2016) & AlphaGo Zero (2017) & AlphaZero (2018) \\
\midrule
Human game data      & yes (supervised policy) & removed & not used \\
Networks             & 3 learned $+$ SL policy  & 1 with two heads & 1 with two heads \\
Handcrafted features & some & removed & removed \\
Random rollouts      & yes (fast policy) & removed & removed \\
Symmetry exploitation & --- & augmentation, averaging & removed \\
Best-network gating  & --- & 55\% gate & removed \\
Games supported      & Go & Go & chess, shogi, Go \\
\bottomrule
\end{tabularx}
\end{table}

\paragraph{AlphaGo (2016).} The system was a pipeline: a supervised policy
network trained on human games, refined by policy-gradient self-play
\eqref{eq:reinforce-game}; a separate value network; a fast rollout policy of
hand-crafted features; and MCTS that combined all four
\citep{silver2016alphago}. It beat Fan Hui 5--0 and Lee Sedol 4--1.

\paragraph{AlphaGo Zero (2017).} Human data, rollouts, hand-crafted features, and
the separate networks were all deleted. One network with two heads, one search,
and a self-play loop remained. After three days of training it beat the
Lee-beating version 100 games to 0; after 40 days it reached an Elo of 5{,}185,
against 4{,}858 for the strongest previous version and 3{,}739 for AlphaGo Lee.
The raw network, with no search at all, measured 3{,}055 \citep{silver2017alphagozero}.
That last number is worth pausing on: it means the network alone, playing a move
per forward pass, was already at the level of a distributed system that used 176
GPUs.

\paragraph{AlphaZero (2018).} The gate, the symmetry exploitation, and the
assumption of binary outcomes were removed, giving one algorithm, one
architecture, and one set of hyperparameters for three games
\citep{silver2018alphazero}. It beat Stockfish after 4 hours, Elmo after 2 hours,
and AlphaGo Lee after 8 hours of training.

\subsection{The algorithm}
\label{sec:algorithm}

\subsubsection{The network}

One network computes both outputs:
\begin{equation}
f_\theta(s) \;=\; (p, v), \qquad
p_a \approx \Prob(a \mid s), \qquad v \approx \E[z \mid s],
\end{equation}
where $s$ is the position, $p$ is a distribution over moves, $v \in [-1,1]$ is
the expected outcome from the side to move, and $z$ is the final outcome of the
game. A single forward pass therefore supplies everything the search needs: a
prior over moves and an evaluation of the position.

\subsubsection{The search}

Each move is chosen after a fixed budget of simulations (1{,}600 in AlphaGo Zero,
800 in AlphaZero; the numbers must not be mixed). One simulation has three
phases, and the difference from classical MCTS is in the second.

\begin{description}
  \item[Select.] Descend from the root, at each node choosing the action of
        \eqref{eq:puct}, until reaching a leaf. A leaf is either terminal or
        never-expanded.
  \item[Expand and evaluate.] If the leaf is terminal, its value is the true game
        outcome, and the network is not called. Otherwise, evaluate it once:
        $(p,v) = f_\theta(s_L)$, store $p$ as the priors of the leaf's outgoing
        edges, and use $v$ as the leaf's value. \emph{There are no rollouts.}
        This is the whole of the replacement of AlphaGo's rollout policy: one
        trained evaluation instead of several hundred random moves.
  \item[Back up.] Propagate $v$ to the root, incrementing visit counts and
        accumulating values along the way. The sign alternates at each ply,
        because the game is zero-sum and the players alternate
        (Section~\ref{sec:tree} takes this up as an implementation matter).
\end{description}

After a move is played, the subtree under the played edge becomes the new root;
the accumulated statistics survive, which is free depth at the next decision.

\subsubsection{Self-play data generation}

Self-play games are generated with the current network guiding the search, with
two exploration mechanisms active at the root.

\emph{Dirichlet noise} perturbs the root priors:
\begin{equation}
P(s_{\mathrm{root}}, a) \;=\; (1-\varepsilon)\, p_a + \varepsilon\, \eta_a,
\qquad \eta \sim \mathrm{Dir}(\alpha), \qquad \varepsilon = 0.25 .
\label{eq:dirichlet}
\end{equation}
AlphaGo Zero used $\alpha = 0.03$ for Go. AlphaZero scaled $\alpha$ in inverse
proportion to the typical number of legal moves: $0.3$ for chess, $0.15$ for
shogi, $0.03$ for Go \citep{silver2018alphazero}. The paper's reason is that the
noise ensures every move may be tried while the search may still overrule bad
moves. A small $\alpha$ produces a spiky noise vector, so a few random moves
receive a large boost per game and different games explore different lines.

\emph{Temperature} controls how the played move is sampled from the visit counts:
\begin{equation}
\pi_a \;\propto\; N(s_{\mathrm{root}}, a)^{1/\tau}.
\label{eq:temperature}
\end{equation}
For the first 30 moves of each AlphaGo Zero game, $\tau = 1$, so play is
proportional to visit counts and openings diversify. Afterwards $\tau \to 0$,
which means playing the most-visited move. The asymmetry is deliberate: early
diversity widens the state distribution that the value head sees, while late
greediness keeps the value target informative, because a game thrown away by
random midgame play teaches the value head very little.

Each position of a finished game is stored as a triple $(s, \pi, z)$, where
$\pi$ is the full visit-count distribution at the root -- not the move that was
played -- and $z$ is the final outcome from the perspective of the player to move
at $s$. AlphaGo Zero additionally applied a random element of the board's
dihedral symmetry group to $s$; AlphaZero removed that step because chess and
shogi have no such symmetry.

AlphaGo Zero also resigned games in which the value estimate stayed below a
threshold, with the false-positive rate held at or below 5\% by forcing 10\% of
games to be played to the end for calibration \citep{silver2017alphagozero}.

\subsubsection{The training step}

Mini-batches are drawn uniformly from a replay window of recent self-play games
(500{,}000 games in AlphaGo Zero; AlphaZero's window size is not published). For
a sample $(s,\pi,z)$ with network output $(p,v)$ the loss is

\begin{equation}
\ell(\theta) \;=\; \underbrace{(z-v)^2\vphantom{\big|}}_{\text{value}}
\;-\; \underbrace{\pi^{\!\top} \log p\vphantom{\big|}}_{\text{policy}}
\;+\; \underbrace{c\,\|\theta\|^2\vphantom{\big|}}_{\text{regularization}},
\qquad c = 10^{-4}.
\label{eq:alphazero-loss}
\end{equation}

The two error terms carry equal weight. AlphaGo Zero justifies this in one
parenthesis: it is reasonable because the rewards are unit-scaled, $r \in
\{-1,+1\}$ \citep{silver2017alphagozero}. Section~\ref{sec:loss} unpacks that
sentence. The optimizer is stochastic gradient descent with momentum $0.9$; the
learning-rate schedules are in Table~\ref{tab:hyper}.

\subsubsection{The outer loop}

AlphaGo Zero ran iterations with a gate: generate games, train a candidate, and
promote the candidate only if it won at least 55\% of 400 evaluation games
against the current best. AlphaZero removed the gate, always training on data
from the latest network. The cost of the gate is compute plus the training steps
spent waiting; the benefit is that a bad network never generates training games.
Removing it was shown to be safe at scale, in part because the replay window
mixes data from many recent checkpoints, so a transient regression dilutes
rather than dominates \citep{tian2019elf,wu2019katago}.

\begin{figure}[t]
\centering
\begin{tikzpicture}[
  node distance=9mm and 14mm,
  box/.style={draw, rounded corners=2pt, align=center, font=\small,
              minimum height=9mm, inner xsep=6pt, fill=black!3},
  arr/.style={-{Latex[length=2mm]}, thick},
  lbl/.style={font=\scriptsize, inner sep=2pt},
]
\node[box] (net) {network $f_\theta$\\$(p,v)$};
\node[box, right=of net] (search) {MCTS\\$\pi = \mathrm{norm}(N)$};
\node[box, right=of search] (games) {self-play games\\$(s,\pi,z)$};
\node[box, below=of games] (buffer) {replay window\\uniform sampling};
\node[box, left=of buffer] (train) {gradient step on\\$\ell(\theta)$};
\draw[arr] (net) -- node[lbl,above]{prior, value} (search);
\draw[arr] (search) -- node[lbl,above]{improved policy} (games);
\draw[arr] (games) -- (buffer);
\draw[arr] (buffer) -- (train);
\draw[arr] (train) -| node[lbl,below,pos=0.25]{updated weights} (net);
\end{tikzpicture}
\caption{The AlphaZero loop. The network is trained to imitate the search; the
improved network then produces a stronger search. Everything else in the system
is engineering around this cycle.}
\label{fig:loop}
\end{figure}

\subsection{Why the loop works}
\label{sec:alphazero-why}

\paragraph{The skeleton is approximate policy iteration.} Section~\ref{sec:apxpi}
already made this identification. The search is the improvement operator, the
value head trained on $z$ is the evaluation operator, and the network is refitted
to both. What makes the loop climb is the asymmetry: the teacher is stronger than
the student, because the teacher is the student plus search.

\paragraph{The value target is the raw game result.} There is no term of the form
$v(s) \leftarrow r + v(s')$. The consequences are worth listing, because they
explain most of the stability of the method.

\emph{Unbiased targets}: $z$ is a sample of the true outcome under the current
improved policy, so it carries no approximation error. \emph{No deadly triad}:
the instability of Section~\ref{sec:actorcritic} requires bootstrapping, and
bootstrapping is absent, so the triad cannot form. \emph{Variance}: a single game
outcome is a noisy estimate of a position's worth; the cure is volume, namely
tens of millions of games and a replay window that averages the noise across each
batch. The team judged unbiased targets worth the variance, in a class of games
where exploration noise must come from the search rather than from dice. KataGo
kept the same choice \citep{wu2019katago}.

\paragraph{The policy target is a distribution.} Training uses $\pi$, the full
visit-count distribution at the root, not the move played. Three arguments
support this. It is an information-dense target: $\pi$ encodes the search's
opinion about every legal move, including how strongly the alternatives were
rejected, while an argmax target keeps one bit of relative preference. This is
the same argument as soft targets in knowledge distillation
\citep{hinton2015distillation}: the relative probabilities of the wrong answers
carry most of the teaching signal. It is a low-variance target, because it
averages hundreds of simulations. And it is by construction the output of the
improvement operator, so imitating it is exactly the imitation step of the
policy-iteration loop.

\paragraph{MCTS averages network errors; minimax amplifies them.} This is the
AlphaZero paper's own argument for why search and deep evaluation fit together,
and it is worth quoting in full:

\begin{quote}
``MCTS averages over these approximation errors, which therefore tend to cancel
out when evaluating a large subtree. In contrast, alpha--beta search computes an
explicit minimax, which propagates the biggest approximation errors to the root
of the subtree'' \citep{silver2018alphazero}.
\end{quote}

A deep network's value errors are confident and hard to predict. Plug such an
evaluator into minimax and the max operator systematically seeks out the largest
errors, because the maximum of a set of noisy estimates is biased upward. MCTS
averages evaluations over many leaf visits, so independent errors partially
cancel. The search is not merely compatible with the network; it is the search
whose error model matches the network's failure mode.

\paragraph{Knowledge replaces computation.} Deleting the rollouts replaced a
cheap, biased, high-variance estimate with one forward pass that returns a
calibrated expectation. The aggregate effect is visible in the search statistics:
in chess AlphaZero searches about 80{,}000 positions per second while Stockfish
searches about 70 million, and in shogi the numbers are 40{,}000 against 35
million \citep{silver2018alphazero}. AlphaZero wins with roughly three orders of
magnitude less search, because each of its evaluations carries learned knowledge
instead of statistics over random play. Breadth is traded for depth of
evaluation, and nothing in the algorithm depends on how fast random rollouts can
be made -- which is also why the algorithm transfers between games.

\paragraph{Self-play is an automatic curriculum.} The opponent is always exactly
as strong as the learner. Early games are between near-random players and are
short and tactically simple; as the network improves, so does the opposition, and
the data always sits at the edge of current ability, which is the regime where
the gradient signal is richest. There is no distribution shift between training
and deployment, because both are self-play. The exploration mechanisms of
Section~\ref{sec:algorithm} exist to prevent the data from collapsing onto a
single deterministic line, which would starve the value head.

\paragraph{Evidence that the loop discovers rather than imitates.} During AlphaGo
Zero training the system rediscovered human joseki, the corner patterns developed
over centuries of play, and then abandoned several of them in favour of variants
it preferred \citep{silver2017alphagozero}. Human knowledge appeared as a
transient phase rather than as a ceiling.

\subsection{Architecture}
\label{sec:architecture}

\paragraph{Input as a stack of planes.} AlphaGo Zero represents a position as a
$19\times19\times17$ image: eight planes for the current player's stones at time
steps $t \dots t-7$, eight for the opponent's, and one constant plane for the
side to move \citep{silver2017alphagozero}. The plane representation is not
incidental. A board is a grid, so the state is an image; the rules are
translation-invariant, which matches weight sharing; they are defined through
adjacencies, which matches local receptive fields; and they are invariant under
rotation and reflection, which matches data augmentation. The paper states this
reasoning explicitly. The history planes supply the temporal context
(repetitions, ko-like situations) that a single snapshot cannot show.

AlphaZero generalized the encoding per game \citep{silver2018alphazero}: keep the
eight-step history, and replace ``stone or no stone'' with one plane per piece
type. Chess needs 6 piece types per side plus 2 repetition planes, times 8
history steps, plus 7 constant planes -- 119 planes per cell. The organizing
principle survived unchanged, and the plans are always oriented to the side to
move, which halves the state space.

\paragraph{Trunk and heads.} The AlphaGo Zero trunk is a convolutional block of
256 filters of size $3\times3$ followed by batch normalization and a rectifier,
then 19 or 39 residual blocks, each of the form
$\mathrm{conv}\,3{\times}3 \to \mathrm{BN} \to \mathrm{ReLU} \to
\mathrm{conv}\,3{\times}3 \to \mathrm{BN} \to (+ \mathrm{skip}) \to \mathrm{ReLU}$

\citep{silver2017alphagozero}. The policy head is a $1\times1$ convolution with
two filters, batch normalization, a rectifier, and a linear layer to one logit
per move; the value head is a $1\times1$ convolution with one filter followed by
linear layers to 256 and then to 1, with a $\tanh$ output.

Two details are load-bearing. The heads are deliberately tiny next to the trunk:
the $1\times1$ convolutions collapse 256 channels to almost nothing before the
expensive linear layers, because the knowledge is expected to live in the shared
body. And the $\tanh$ bounds $v$ to the range of $z$, which keeps the squared
error well scaled and prevents occasional large logits from producing large
gradients.

\paragraph{Why the residual tower, and why one tower for two heads.} The
architecture is the ResNet design \citep{he2016resnet} transplanted from image
classification, and the paper says the hyperparameters were chosen to match the
state of the art in image recognition. Two arguments matter.

Residual connections make very deep networks trainable at all: without them
gradients vanish and training accuracy degrades with depth. And the choice was
measured rather than assumed. AlphaGo Zero trained four architectures on the same
fixed dataset of self-play games \citep{silver2017alphagozero}: one residual
tower with both heads; two residual towers, one per task; one non-residual
convolutional tower with both heads; and two non-residual towers. The first won.
The residual tower accounts for more than 600 Elo, and combining policy and value
into one network accounts for roughly another 600 Elo. The paper reports that
the combination slightly reduced move-prediction accuracy while reducing value
error and raising playing strength, and explains it as regularization to a common
representation that serves multiple use cases.

Two mechanisms produce that effect. Computationally, MCTS needs both outputs at
every leaf, so two separate towers would cost two forward passes per leaf, which
is the dominant cost of self-play. Statistically, the value task forces the trunk
to encode features that predict the outcome while the policy task forces it to
encode features that discriminate good moves from bad; neither task alone needs
the full concept set of the other, but a shared trunk must serve both, so it
learns the union. Each objective trains the other's features. The small drop in
raw move-prediction accuracy is the signature of the trade: a little policy fit
is given up for a lot of position understanding, and playing strength follows
position understanding.

\paragraph{Action encoding.} Go uses a flat distribution over the board plus a
pass move, which does not generalize to games with structured moves. AlphaZero
encodes chess moves as an $8\times8\times73$ stack -- 56 queen-move planes, 8
knight-move planes, 9 underpromotion planes -- giving 4{,}672 moves, and shogi as
$9\times9\times139$ \citep{silver2018alphazero}. Keeping the output spatial, so
that a move means ``from this square, in this direction, by this much'',
preserves the same translation symmetry that helps the input. The price is a
fixed maximal move set, with illegal moves masked and renormalized at search
time. For a game whose moves are simply ``place a stone on an empty cell'' the
flat encoding is already spatial: one output per cell, masked to the legal cells.

\subsection{The loss, term by term -- and why it is not a policy gradient}
\label{sec:loss}

Equation \eqref{eq:alphazero-loss} has three terms with three jobs.

\paragraph{$(z-v)^2$: policy evaluation.} Fit the value head to the true outcome
of games played by the improved policy. The paper treats $v$ as an expectation,
$v \approx \E[z \mid s]$, and mean squared error is the natural fit for an
expectation of a bounded scalar. It preserves the unit-scale argument that
justifies equal weighting, and it accommodates the draw value $0$ without a
reparametrization.

\paragraph{$-\pi^{\!\top}\log p$: policy improvement, written as distillation.}
Cross-entropy between the search distribution and the network policy. Its
minimizer is $p = \pi$: the network is asked to reproduce the search's output for
every move, not only for the move played. This is the term that makes the loop
climb, because $\pi$ is the output of an improvement operator, so fitting $\pi$
pulls the network above its own raw-policy level at every iteration, with
supervision generated on the fly.

\paragraph{$c\|\theta\|^2$: smoothness, with a system-level purpose.} In this
system the regularizer does more than prevent overfitting. MCTS trusts the
network: the prior steers selection and the value steers backup. A network that
overfits the replay window becomes locally erratic -- priors spike on memorized
positions, values swing between neighbouring positions -- and the search built on
top of it degrades, which degrades the next batch of data. Weight decay keeps the
function smooth, which keeps the search meaningful, which keeps the data stream
healthy. The regularizer protects the loop, not only the fit.

\paragraph{What is deliberately absent.} The loss is also defined by what it does
not contain: no bootstrapped term, so targets stay unbiased; no entropy bonus,
because exploration is the search's job
(Section~\ref{sec:alphaepsilon} explains how that separation is kept); no
importance weights, because the replay window is fresh enough for the
policy-iteration framing; and \emph{no policy-gradient term at all}. The
reinforcement-learning signal enters through the data -- the search results --
and not through the gradient estimator.

This last point is the one that is most often misstated. AlphaZero does not
estimate $\nabla_\theta J$ by any method resembling REINFORCE. It runs policy
iteration and supervises the pieces. The policy-gradient layer of
Section~\ref{sec:pgt} is the historical ancestor: it is what trained AlphaGo's
reinforcement-learning policy network in 2016 \eqref{eq:reinforce-game}. By 2017
the explicit gradient step was gone, replaced by distillation from search. The
object being optimized is still the expected outcome of a parameterized policy,
and the convergence intuition still comes from the policy gradient theorem, but
the mechanism is a supervised-fit-to-improved-targets loop. That difference is
why the method trains stably at scales where direct policy-gradient methods are
fragile.

\subsection{Hyperparameters}
\label{sec:hyper}

Table~\ref{tab:hyper} collects the published numbers. It is short on purpose:
the papers specify the algorithm in full and the constants barely at all.

\begin{table}[t]
\centering
\small
\caption{Published training quantities. Where a value is absent, the primary
source does not publish it, and the substitute must be treated as an experiment.
Simulations per move are 1{,}600 in AlphaGo Zero and 800 in AlphaZero; the two
must not be mixed.}
\label{tab:hyper}
\begin{tabularx}{\textwidth}{@{}l X X@{}}
\toprule
Quantity & AlphaGo Zero & AlphaZero \\
\midrule
Residual blocks & 20 (3-day), 40 (40-day) & 19 \\
Channels & 256 & 256 \\
Mini-batch & 2{,}048 & 4{,}096 \\
Training steps & 700k (3-day), 3.1M (40-day) & 700k (mini-batches) \\
Self-play games & approx.\ 4.9M & 44M chess, 24M shogi, 21M Go \\
Learning rate & $10^{-2} \to 10^{-3} \to 10^{-4}$ &
$0.2 \to 0.02 \to 0.002 \to 0.0002$ \\
Optimizer & SGD, momentum 0.9 & SGD, momentum 0.9 \\
Weight decay & $10^{-4}$ (inside the loss) & $10^{-4}$ (inside the loss) \\
Replay window & 500{,}000 games & not published \\
$c_{\mathrm{puct}}$ & not published & not published (1.5 is ELF's value) \\
Dirichlet $\varepsilon$ & 0.25 & 0.25 \\
Dirichlet $\alpha$ & 0.03 (Go) & 0.3 chess, 0.15 shogi, 0.03 Go \\
Temperature & $\tau = 1$ for 30 moves, then $\to 0$ & same form \\
Gating & 55\% over 400 games & removed \\
Hardware & 64 GPU workers $+$ 19 CPU servers & 5{,}000 TPUv1 self-play, 64 TPUv2 train \\
\bottomrule
\end{tabularx}
\end{table}

\subsection{The longer history: a decade of deletions}
\label{sec:history}

The design looks inevitable in hindsight, which it was not. It is the residue of
roughly ten years of removing one crutch at a time \citep{silver2009thesis}.
In 2007, reinforcement learning over hand-crafted local pattern features produced
weak amateur play in Go: local patterns alone do not scale. In 2007--2008,
combining UCT with learned pattern priors produced the first dan-level Go
program, which established that learned knowledge inside tree search multiplies
strength \citep{gelly2007rave}. In 2015, deep Q-networks showed that a single
convolutional network trained from pixels could reach human level on a suite of
Atari games \citep{mnih2015dqn}; the result demonstrated that deep function
approximation is stable enough for reinforcement learning, and it is what made
the following step plausible. It also showed the limits of the value-only
approach: DQN learns a value over actions and acts greedily on it, with no policy
network and no search, so it cannot represent a distribution over good moves and
cannot improve itself by lookahead. In 2016, deep networks supplied that policy
as well, and MCTS combined policy, value, rollouts, and human data. In 2017 the
human data, the rollouts, and the separate networks were removed. In 2018 the
gating, the symmetry exploitation, and the Go-specific structure were removed. In
2020 MuZero removed the rules themselves, learning a latent dynamics model so
that the same loop also masters Atari \citep{schrittwieser2020muzero}.

Read as a sequence, every step deleted the most human part that remained, and
every deletion either held or raised playing strength.

\paragraph{The ancestor.} Tesauro's TD-Gammon (1992--1995) trained a value
network purely by self-play temporal-difference learning and reached strong
master level in backgammon, the first demonstration that self-play with function
approximation can produce world-class play \citep{tesauro1995tdgammon}. It differs
from AlphaZero in three ways that are instructive. Backgammon has dice, which
supply exploration noise for free; AlphaZero must manufacture that noise, which
is why Dirichlet noise and temperature exist. TD-Gammon played from the raw
network with shallow lookahead, so there was no search in the loop. And it
bootstrapped, which AlphaZero does not. The paper's judgement was that unbiased
targets were worth their variance in deterministic games; TD-Gammon is the
counterexample that shows bootstrapping can work when the environment itself
supplies the noise.
